\documentclass{article}
\usepackage{amsmath}

\title{Study Roadmap: Chapters from \textit{College Algebra} by M. Richardson}
\author{Your Name}
\date{}

\begin{document}

\maketitle

\section*{Introduction}
This document outlines the chapters from \textit{College Algebra} by M. Richardson that are most relevant to your goals of strengthening your foundation in algebra, complex numbers, and preparing for FFT and signal processing. The chapters are divided into high priority, secondary priority, and those less relevant for now.

\section*{1. High Priority Chapters}

\subsection*{Chapter 4: Functions and Graphs}
Understanding functions and their graphical representations is fundamental for transformations in linear algebra and FFT. You’ll also encounter this concept frequently in signal processing.

\subsection*{Chapter 5: Elementary Operations of Polynomials}
Polynomials are essential for understanding signal breakdown in FFT and many algorithms in linear algebra. This chapter will give you the base you'll need.

\subsection*{Chapter 9: Integral and Fractional Exponents}
This chapter will help with transformations involving exponents, which are crucial when dealing with complex numbers and FFT.

\subsection*{Chapter 11: Quadratic Equations and Functions}
Quadratics will help strengthen your ability to solve for roots, which ties into signal processing and the mathematics of transformations.

\subsection*{Chapter 14: Complex Numbers}
This is critical for FFT and working with signals in the frequency domain. You’ll want to master this section.

\subsection*{Chapter 15: Theory of Equations}
This chapter ties into how polynomials and equations behave, which will help when understanding systems of equations in FFT.

\subsection*{Chapter 16: Determinants and Elimination Theory}
This is foundational for linear algebra and matrix operations, which are key for FFT and understanding more advanced mathematical transformations.

\subsection*{Chapter 25: Euclidean Algorithm}
The Euclidean algorithm has applications in solving systems of equations, which can be beneficial for linear algebra.

\subsection*{Chapter 27: Infinite Series}
Since FFT and many other algorithms rely on series and approximations, this chapter could give you valuable insights.

\section*{2. Secondary Priority Chapters (Helpful but Not Critical)}

\subsection*{Chapter 7: Elementary Operations with Fractional Expressions}
Understanding these operations will help with algebraic manipulation, though it may not be critical for signal processing.

\subsection*{Chapter 10: Radicals}
This could be useful for general algebraic fluency but is not essential for FFT.

\subsection*{Chapter 12: Systems of Equations in Two Unknowns involving Quadratics}
Systems of equations come up in linear algebra, and while this chapter focuses on quadratics, it can help reinforce broader concepts.

\subsection*{Chapter 18: Probability}
Probability has its applications in signal processing and algorithms, though it’s more of a secondary topic for FFT.

\section*{3. Chapters to Skip for Now (Less Relevant for FFT)}

\subsection*{Chapter 1: Number System}
This is likely a review, not essential for your current focus.

\subsection*{Chapter 2: What is Algebra}
Introductory material you can skip if you’re already comfortable with algebra.

\subsection*{Chapter 24: Mathematics of Investment}
Not relevant for your goals.

\end{document}
